\subsection*{Size 27 and the obstruction at the next size}
At size 27 the exact minimum is $m(27)=72$, and three modulus sets attain it. To prove the lower bound, note that a 27-set through 71 could use at most two entries outside the cap-60 pool, namely 63 and 70. Every 25-subset of that pool is coverable. At 72 the only new modulus is 72, so a candidate needs all of $63,70,72$. Its 24 old moduli cannot cover, and removing 72 cannot leave any proper covering subfamily. This reduces the computation to the surviving cap-70 candidate lists, including those which previously failed to cover. A lone member of the pair $16,48$ is again redundant. There are 22 candidate 27-sets after these exclusions. The same exact ternary-fibre procedure allows the additional 9-divisible row 72 (projected modulus 24). It returns exactly three passing sets. In a concise notation they are
\[
 (P_{60}\setminus\{2,3,e\})\cup\{63,70,72\},
 \qquad e\in\{12,18,24\},
\]
where $P_{60}$ is the entire 27-element candidate pool through cap 60. A full-period residue witness with $e=24$ accompanies the source. A separate binary-digit projection rechecks its cover and rules out each of its 27 deletions. Lean checks coverage over all integers, cardinality 27, pairwise distinctness and maximum 72.

To enumerate residue assignments efficiently one may independently permute the child labels at every node of every prime-power residue tree. Such permutations map a class modulo $d$ to another class modulo $d$ simultaneously for all $d$ dividing the ambient period. For an ordered list of rows, at each prefix name children in order of first occurrence. Trying any already used child and the first unused child enumerates all assignment orbits. In the ternary projection use this normalization for the 9-divisible rows, leaving the full family of 7-masks available. This preserves the quantifiers over assignments and all their deletions.

We also find $m(28)\ge84$. At size 28 through cap 83 all four available additions $63,70,72,80$ to the cap-60 pool are necessary. There are only 35 possibilities for the other 24 rows, since each must form a noncoverable set. The new feature is that the highest binary digit can now have three rows, $16,48,80$. With two rows the usual contraction applies. With all three, their maximal allowable projected hole set is a translate of $(0\bmod24)\cup(0\bmod40)$. Indeed one of the two binary-digit values receives a singleton row. If the singleton projects to modulus 24 or 40, the allowed holes lie in one such class. If it projects to modulus 8 then the other side can supply one class each modulo 24 and 40, restricted to that 8-class. If both are nonempty they agree modulo $\gcd(24,40)=8$, so a common translation puts both at zero. Conversely these two classes together are simultaneously attainable by putting the 8-class on one side and the 24- and 40-classes on the other. Thus the composite row on the 120-cycle has precisely the 120 possible translates of that union. The seven- and ternary-digit method applies unchanged to all the remaining rows. None of the 35 candidates is irreducible. The general lower proof is by projection and exhausts every residue reassignment; the supplied code uses integer bit masks for the finite searches.
